\subsection{The temporal validator}
\label{sec:validator}

The validator is the component this paper measures most closely. Its
mechanism is conventional and not claimed as novel (\cref{sec:related}); the
contribution is the account, given here with controls, of what it buys, what
it costs, and where a simpler baseline would do as well.

\subsubsection{Rationale}

Detection is a commodity: datasets are public, models are published, and
fine-tuning one is a weekend's work. The unsolved, expensive and
operationally decisive problem is the false alarm, which is why commercial
operators surround good models with expensive human staffing.

The validator sits between the detectors and the alert channel. It receives
every raw detection and decides which become events. Detections it rejects are
not discarded silently: each is annotated with the name of the level that
rejected it, so suppression is auditable in evaluation and visible in the
operator interface.

\subsubsection{The cascade}

\begin{definition}[Cascade]
\label{def:cascade}
Let $x$ be a detection with box $b(x)$, confidence $s(x)$ and frame timestamp
$t$. The validator applies three predicates in order, cheapest first:
\begin{align}
  L_1(x) &\iff s(x) \ge \theta
    &&\text{confidence} \label{eq:l1}\\
  L_2(x) &\iff h(\kappa(x)) \ge N
    &&\text{persistence} \label{eq:l2}\\
  L_3(x) &\iff \neg\,\exists\,(b', t') \in R_t :
                \iou(b(x), b') \ge \tau_c
    &&\text{cooldown} \label{eq:l3}
\end{align}
where $\kappa(x)$ is the track $x$ was associated with, $h(\cdot)$ is that
track's hit count, and $R_t = \{(b', t') : t - t' \le \Delta\}$ is the set of
recently confirmed boxes. A detection becomes an event iff
$\bigwedge_{i=1}^{3} L_i(x)$.
\end{definition}

Association is by greedy $\iou$ overlap: a detection joins the existing track
maximising $\iou$ above a threshold $\tau_t$, or starts a new one, and a track
is dropped after $M$ consecutive misses. Every level is independently
switchable. That is not a convenience — it is what allows the ablation of
\cref{sec:ablation} to be produced by the same code that runs in production,
and not by a separate analysis script that might diverge from it.

\subsubsection{Cascade design and level ablation}
\label{sec:removed-levels}

The cascade did not start with three levels. Two more were implemented,
measured under the same harnesses, and removed on the evidence; both removals
are guarded by regression tests, and both are reported here because a cascade
that only ever grows is a cascade nobody is checking.

\begin{recorded}{Size plausibility: harmful on the worst cases}
A size-plausibility level rejected detections whose area fraction or aspect
ratio was implausible. Measured, it rejected 0\,\% of detections on wildfire
and 0\,\% on road incidents. When flood was added it proved actively harmful:
real ATLANTIS flood masks cover a median 32\,\% of the frame and severe cases
exceed 60\,\%, so a plausible maximum-area default silently discarded the
worst floods — the more severe the disaster, the more likely the alert
disappeared. A level that earns nothing on two hazards and inverts the safety
property on a third does not belong in a life-safety cascade. Removing it
changed no measured result, which is itself the evidence it was doing nothing.
\end{recorded}

\begin{recorded}{Colour prior: fog is exactly as grey as smoke}
A colour level checked the fraction of a box's pixels inside a per-hazard HSV
region. Every shipped configuration disabled it, which meant an earlier
draft of this paper defined a level that had never run in a reported
experiment — the same situation that triggered the size measurement. The only
physically sensible prior for wildfire smoke is achromaticity, and as first
parameterised the level could not even express that (it could demand
\emph{minimum} saturation, never a maximum); the bound was added so the
measurement could run at all. Swept over five cutoffs on \ColourStillsImages{}
pyro-sdis stills, the prior bought false alarms only by discarding real
smoke — recall \ColourBaselineRecall{} $\to$ \ColourStrictRecall{} to move
FAR \ColourBaselineFAR{} $\to$ \ColourStrictFAR{} at the strictest cutoff,
about five points of recall per ten of FAR at every point — because fog, the
curated hard negative, shares the one property being tested. Inside the
temporal cascade it improved nothing at any cutoff, and strict cutoffs lost
an ignition sequence outright (\ColourStrictFires{} of \FiresTotal). Removed
(\texttt{eval/run\_colour\_eval.py}; the harness reimplements the prior
locally so the experiment outlives the level it deleted).
\end{recorded}

\subsubsection{Decomposing suppression into filtering and deduplication}
\label{sec:decomposition}

The validator rejects the large majority of what the detectors propose, and it
is tempting to report that as a single suppression figure. That figure would
be misleading, because it combines two operations supporting entirely
different claims.

\begin{definition}[Suppression decomposition]
\label{def:decomposition}
Let $P$ be the number of proposals and $E$ the number of confirmed events.
Write $\Phi$ for the proposals rejected by $L_1$--$L_2$ and $D$ for those
rejected by $L_3$, and define
\begin{equation}
  S_{\text{total}} = 1 - \frac{E}{P},
  \qquad
  S_{\text{filter}} = \frac{\Phi}{P},
  \qquad
  S_{\text{dedup}} = \frac{D}{P}.
  \label{eq:decomposition-terms}
\end{equation}
Since every proposal is either confirmed or rejected by exactly one level,
\begin{equation}
  S_{\text{total}} \;=\; S_{\text{filter}} \;+\; S_{\text{dedup}}.
  \label{eq:decomposition}
\end{equation}
\end{definition}

Only the first term is false-alarm reduction, and only the first term supports
the claim this project makes. The second is the system declining to
re-announce an event it has already reported: nothing is filtered out there,
an alert is simply not sent twice.

Measured across the four demonstration clips, \SuppressionProposed{} proposals
produced \SuppressionConfirmed{} confirmed events. As one number that is
\SuppressionTotal\,\% suppression. Split by \cref{eq:decomposition} it is
\SuppressionFiltered\,\% filtering and \SuppressionDeduplicated\,\%
deduplication — and the split varies enormously by hazard.

\begin{table}[htbp]
\centering\small
\caption{The decomposition of \cref{eq:decomposition} per hazard. The
proportions invert between fire and building access, which a single
suppression figure would conceal entirely. \measured{} —
\texttt{docs/figures/suppression.json}.}
\label{tab:decomposition}
\begin{tabular}{l S[table-format=3.0] S[table-format=3.0]
                  S[table-format=2.1] S[table-format=2.1]}
\toprule
Hazard & {Proposed} & {Confirmed} & {Filtering \%} & {Dedup. \%} \\
\midrule
Fire             & \SplitFireProposed           & \SplitFireConfirmed           & \SplitFireFiltered           & \SplitFireDeduplicated \\
Flood            & \SplitFloodProposed          & \SplitFloodConfirmed          & \SplitFloodFiltered          & \SplitFloodDeduplicated \\
People in water  & \SplitDrowningProposed       & \SplitDrowningConfirmed       & \SplitDrowningFiltered       & \SplitDrowningDeduplicated \\
Building access  & \SplitBuildingAccessProposed & \SplitBuildingAccessConfirmed & \SplitBuildingAccessFiltered & \SplitBuildingAccessDeduplicated \\
\bottomrule
\end{tabular}
\end{table}

\begin{figure}[htbp]
\centering
\includegraphics[width=0.9\textwidth]{suppression_split.pdf}
\caption{What the validator does with each hazard's proposals. Fire's headline
suppression is almost entirely the system declining to re-announce one fire it
had already found; building access's is almost entirely filtering. Quoting a
single figure would overstate the false-alarm contribution by a factor of
eight on the fire clip.}
\label{fig:suppression}
\end{figure}

\begin{recorded}{This project has been caught by exactly this before}
During the first ablation, recall appeared to collapse from $0.925$ to $0.125$
and read as the validator destroying the detector. It was cooldown correctly
suppressing repeat alerts about a fire already found — the validator was
locating \emph{more} fires, not fewer. Cooldown was removed from the
evaluation harness for that reason. The runtime pipeline later reintroduced
the same conflation in its live statistics, and now reports the two quantities
separately everywhere they appear, including on the recorded demonstrations
and in the operator interface.
\end{recorded}

\subsubsection{False-alarm reduction}
\label{sec:ablation}

\begin{table}[htbp]
\centering\small
\caption{The validator against the raw detector on
\AblationNegativeSequences{} negative sequences and \AblationFireSequences{}
ignition sequences from HPWREN FIgLib. \measured{} —
\texttt{eval/run\_ablation.py}.}
\label{tab:ablation}
\begin{tabular}{l S[table-format=1.3] S[table-format=2.0] c
                  S[table-format=1.3] S[table-format=1.1]}
\toprule
Configuration & {FAR} & {Reduction \%} & Fires found & {Frame recall}
  & {Min.\ to detect} \\
\midrule
Raw detector, no validator      & \RawFAR       & {—}            & \RawFiresFound/\FiresTotal & \RawFireRecall       & \RawMinutesToDetect \\
Full cascade, persistence $= 3$ & \ValidatedFAR & \FARReduction  & \FiresFound/\FiresTotal    & \ValidatedFireRecall & \MedianMinutesToDetect \\
\bottomrule
\end{tabular}
\end{table}

At persistence 3 the false-alarm rate falls by \FARReduction\,\% and every
fire is still found: \FiresFound{} of \FiresTotal{}, the same as the raw
detector. Sequence-level detection is what matters operationally, and it is
unchanged. A ByteTrack-style second association stage~\citep{zhang2022bytetrack}
— below-floor detections keeping tracks alive without being allowed to alarm
— was also implemented and measured: it changed nothing on either the sparse
ablation or a dense \SI{8}{fps} sequence, because the tracker's miss
tolerance already bridges the gaps it exists to bridge. It ships disabled, as
a documented null result.

Two costs are paid for that, and both sit in \cref{tab:ablation}, not
in a footnote. The median time to first alert rises from
\RawMinutesToDetect{} to \MedianMinutesToDetect{} minutes, and frame-level
recall falls from \RawFireRecall{} to \ValidatedFireRecall{} — the frames
spent accumulating persistence are frames on which the system stays silent
about a fire it can already see. Reporting only the false-alarm reduction
would hide both. A filter that drove false alarms to zero by refusing to alert
at all would score perfectly on the one metric and be useless.

\subsubsection{Comparison with confidence thresholding}
\label{sec:threshold-matched}

\Cref{tab:ablation} compares the cascade against the raw detector \emph{at
the same confidence floor}. That leaves open the objection any careful
reader raises next: forget temporal validation — raise the floor until the
false-alarm rate matches, and see what it costs. Most system papers leave
this control to the reviewer. We ran it, and it partially deflates our own
headline, so it is reported in full.

Sweeping the raw detector's floor over the same cached detections, a floor
of \MatchedFloor{} —
selected \emph{on these very negatives}, so an oracle baseline — reaches the
cascade's false-alarm rate of \CascadeFAR{} while keeping frame recall
\MatchedRecall{} against the cascade's \CascadeRecall{}, all
\MatchedFires{}/\FiresTotal{} fires, and a median time-to-alert of
\MatchedMinutes{} minute against the cascade's \MedianMinutesToDetect. A
floor of \ZeroFARFloor{} drives the false-alarm rate to zero on this set
outright, at frame recall \ZeroFARRecall. \Cref{fig:frontier} shows the
frontier; the cascade's operating point sits below it.

\begin{figure}[htbp]
\centering
\includegraphics[width=0.9\textwidth]{threshold_frontier.pdf}
\caption{The control experiment
(\texttt{eval/run\_threshold\_matched.py}). Each blue point is the raw
detector at one confidence floor, on identical cached detections; the orange
point is the cascade at persistence 3. On this negative set an oracle-tuned floor
dominates the cascade — the finding this paper reports against itself, and
the reason the claim in \cref{sec:scoped-claim} is scoped the way it is.}
\label{fig:frontier}
\end{figure}

\paragraph{What the control does and does not establish}
\label{sec:scoped-claim}

The control establishes that \textbf{per-frame false-alarm reduction on this
negative set is not the cascade's contribution}: a tuned threshold does it
more cheaply. Three things bound how far that conclusion travels, each a
fact, none a defence:

\begin{itemize}
  \item \textbf{The matched floor is an oracle.} The sweep selects
        $\theta = \MatchedFloor$ using the same 30 sequences it is scored on;
        the cascade runs the detector authors' published default ($0.20$)
        untouched. A deployed threshold must be chosen on some cameras and run
        on others, so the honest comparison is not oracle-against-cascade but
        transferred-threshold-against-cascade. We run that comparison in
        \cref{sec:threshold-transfer-exp}, and it reverses the control's
        verdict.
  \item \textbf{No floor survives the hard negatives.} On the curated
        \texttt{pyro-sdis} negatives (\cref{sec:fire}), raising the floor to
        $0.30$ still leaves a false-alarm rate of $0.648$ while recall is
        already falling — fog and dust are detected \emph{with high
        confidence}. Whether persistence helps there is untested, because no
        temporal hard-negative set exists in any release we know of; building
        one is listed in \cref{sec:outstanding}.
  \item \textbf{Deduplication is not achievable by any threshold.}
        \Cref{tab:decomposition} shows most of what the validator does on
        several hazards is withholding repeats of an event already reported —
        an operational function a per-frame threshold cannot express at all.
\end{itemize}

\paragraph{The transfer experiment: the objection, answered}
\label{sec:threshold-transfer-exp}

The control's central caveat — that its winning floor is an oracle — is
testable, so we test it. The \TransferCameras{} FIgLib negative cameras are
split into two disjoint folds; on one fold the floor is tuned to the target
false-alarm rate, that fixed floor is applied to the \emph{other} fold, and
the result is compared against the cascade running its untuned published
floor on the same held-out cameras (\texttt{eval/run\_threshold\_transfer.py}).
The split is deterministic — cameras sorted and alternated — so no seed has to
be trusted, and both directions are reported.

The oracle does not survive transfer. Tuned to a false-alarm rate of
\TransferTrainFAR{} on one fold (floor \TransferTunedFloor), it drifts to
\TransferTestFAR{} on the cameras it has not seen — more than three times the
target — while the cascade at the published floor holds \TransferCascadeFAR{}
on the same cameras. Tuned the other way (floor \TransferBATunedFloor) it
does hold the false-alarm rate on transfer, but only by \emph{losing a fire}:
recall \TransferBARecall{} against the cascade's \TransferBACascadeRecall.
Averaged over both directions the transferred threshold sits at
\TransferMeanThresholdFAR{} against the cascade's \TransferMeanCascadeFAR.

This reverses the control of \cref{sec:threshold-matched}. A confidence floor
beats the cascade only with oracular access to the negatives it will face; a
floor chosen the way a deployment must choose it either lets false alarms
drift up or sacrifices a detection, whereas the cascade's persistence — being
a property of the tracks, not of any camera's confidence calibration — holds
both false-alarm rate and recall across cameras it has never seen. The
camera-independence of persistence, asserted by an earlier draft, is here a
measured advantage.

The honest statement of the validator's value is therefore this, and it is
now supported by the experiment above rather than by assertion: \emph{a
single untuned mechanism that holds its false-alarm rate and its recall
across unseen cameras where a tuned threshold holds at most one of the two,
provides deduplication no threshold can express, and transfers across hazards
by configuration} — at a measured cost in latency (\cref{sec:latency}) that
is predictable in closed form.
Where a deployment can tune a floor on representative negatives and does not
need deduplication, thresholding alone is the better choice, and this paper
says so.

\subsubsection{Latency cost}
\label{sec:latency}

The latency is not an empirical surprise. Persistence requires $N$ associated
sightings before confirming, so for a camera sampling at rate $f$ the
additional delay relative to a single-frame detector is
\begin{equation}
  \mathbb{E}[\Delta t] \;=\; \frac{N - 1}{f},
  \label{eq:latency}
\end{equation}
assuming the track is not missed. HPWREN's FIgLib cameras sample once per
minute, so at $N = 3$ \cref{eq:latency} predicts exactly
$(3-1)/1 = \SI{2}{\minute}$ of extra delay. The measured medians are
\RawMinutesToDetect{} minutes raw and \MedianMinutesToDetect{} validated, a
difference of \AddedMinutesToDetect{} minutes — the predicted value. The
identity matters more than the agreement: it says the latency cost is a
property of the \emph{camera's sampling rate}, not of the method, so the same
cascade on a camera sampling at \SI{1}{\hertz} costs two seconds. A deployment
choosing $N$ is choosing a point on \cref{eq:latency}, and can be told what it
will cost before any data is collected.

\subsubsection{Cross-hazard transfer}

PyroNear's published training manifest shows that they already perform
sequential confirmation for wildfire, with a grid search over a minimum frame
count. Temporal confirmation for fire is therefore not novel, and we say so.
The claim VIGÍA makes is narrower and still defensible: \emph{one} validator,
configured rather than specialised, working across hazards whose detectors
were built by different people, on different data, for phenomena with
different physics.

\begin{table}[htbp]
\centering\small
\caption{The same validator on two hazards with different physics: a wildfire
plume develops over tens of minutes and is nearly static between frames; a
collision resolves in under a second with objects moving fast across the
frame. The only difference between the two runs is a configuration object.
\measured{} — \texttt{eval/run\_multihazard.py}.}
\label{tab:transfer}
\begin{tabular}{l l S[table-format=3.0] S[table-format=3.0]
                  S[table-format=2.0] S[table-format=2.1]}
\toprule
Hazard & Source & {Frames} & {Raw det.} & {Confirmed} & {Suppression \%} \\
\midrule
Wildfire smoke  & FIgLib, 1 frame/min & \FireFrames    & \FireRawDetections    & \FireConfirmedEvents    & \FireSuppression \\
Road incidents  & Dashcam, 2 fps      & \TrafficFrames & \TrafficRawDetections & \TrafficConfirmedEvents & \TrafficSuppression \\
\bottomrule
\end{tabular}
\end{table}

\subsubsection{Scope of the evidence}
\label{sec:evidence-reach}

The scope of what has and has not been measured deserves an explicit
statement, because the two are easy to conflate. Three distinct kinds of
false-alarm evidence appear in this paper, on different numbers of hazards:

\begin{enumerate}
  \item \textbf{Per-frame detector false-alarm rate}, measured on curated or
        natural negatives, is reported for \emph{three} hazards: fire
        (\FireImageFAR{} on \FireNegatives{} negatives, \cref{sec:fire}),
        people in water (\SwimmerImageFAR, \cref{sec:piw}) and building
        access (\BuildingImageFAR, \cref{sec:building}). This establishes
        that the false-alarm problem the paper motivates is real beyond the
        one detector used for the headline.
  \item \textbf{The temporal validator's false-alarm \emph{reduction}} — the
        cascade against the raw detector on negative sequences, with the
        transfer control — is demonstrated on \emph{one} hazard, fire,
        because it is the only hazard for which negative temporal sequences
        exist in released form. The flood trend detector's false-alarm
        behaviour is measured separately, on stable water with realistic
        segmentation noise (\cref{sec:ph-calibration}), which is a second
        hazard's false-alarm measurement by a different mechanism.
  \item \textbf{Configuration-only transfer} of the cascade across hazards
        with different physics is shown in \cref{tab:transfer}, but that
        table measures suppression, not accuracy against held-out labels.
\end{enumerate}

The honest consequence is that the central causal claim — \emph{temporal
validation reduces the false-alarm rate} — is evidenced directly on fire and
by mechanism-analogue on flood, and that extending it to a second detector
with negative sequences is the single most valuable experiment this work has
not been able to run (\cref{sec:outstanding}). The paper's title says
multi-hazard because the \emph{system} spans four hazards and the validator
transfers to them by configuration; it does not claim the reduction is
independently proven on all four.

\begin{recorded}{Enforced, not merely asserted}
\texttt{tests/test\_validator.py} contains
\texttt{test\_validator\_contains\_no\_hazard\_specific\_code}, which parses
the validator's source with Python's AST module and fails the build if any
comparison against a specific hazard appears. A future contributor fixing a
wildfire bug with a hazard-specific branch cannot silently invalidate the
central claim of this document. The suite passes 9 of 9.
\end{recorded}
